% ===== 原文 pp.24–27（PDF p.34–37） =====
% [原文 p.24]
方程的主象征只依赖于最高阶项的形式，而该形式在 $M$ 的坐标变换下不变。在某点的法坐标下，并进一步利用方程在靶空间坐标变换下的不变性，可以推出这些微分方程的椭圆性。

\subsection{局部可观测量}

我们简要讨论 4d A-模型的局部可观测量谱，即 BRST-闭的局部算子模去 BRST-正合的局部算子。为了写出场的 0-形式泛函，人们可以选择用度量 $h_{\mu\nu}$ 来缩并 1-形式与 2-形式场的指标。然而，如此得到的可观测量自动是 BRST-正合的。因此，在局部可观测量的每个 BRST 上同调类中，都可以选出一个代表，它是 0-形式场 $\phi^I$ 与 $\zeta^I$ 的（BRST-闭的）泛函 $\cO$。

% [Tutor Note: p.24 局部可观测量的鬼数必须与形式次数一致，可用 (3.1) 核对：$\phi,B,\tilde H$ 的鬼数为 $0$，$\zeta,\tilde\psi$ 为 $1$，$\tilde\chi$ 为 $-1$，而 $\delta_Q$ 一律把鬼数加 $1$（$\delta_Q\phi=\zeta$、$\delta_QB=\tilde\psi$、$\delta_Q\tilde\chi=\tilde H$ 均如此）；因此 $\cO$ 的鬼数 $k$ 恰是 $W_{\cO}$ 的形式次数，$\delta_Q$ 对应外微分 $\dd$ 才可能成立。若某式的鬼数不匹配，则须怀疑该式有误。]

给定这样一个鬼数为 $k$ 的泛函 $\cO$，有一个显然的办法写出 $X$ 上相应的微分 $k$-形式 $W_{\cO}$，即
\[
  \cO=\frac{1}{k!}f_{I_1\ldots I_k}(\phi)\,\zeta^{I_1}\cdots\zeta^{I_k}
  \quad\longleftrightarrow\quad
  W_{\cO}=\frac{1}{k!}f_{I_1\ldots I_k}(\phi)\,\dd\phi^{I_1}\wedge\cdots\wedge\dd\phi^{I_k}.
\]
由 BRST 变分 (3.1) 可见，场的 BRST 变分对应于 $X$ 上形式的外微分，即
\[
  W_{\delta_Q\cO}=\dd W_{\cO}.
\]
因此，局部可观测量的 BRST 上同调类对应于 de Rham 上同调类
\[
  H^\bullet_{\text{de Rham}}(X).
\]

% [Tutor Note: p.24 $W_{\delta_Q\cO}=\dd W_{\cO}$ 的实质是 $\delta_Q\phi^I=\zeta^I$ 且 $\zeta^I$ 为费米子（反对易）。核对方法：对 $W_{\cO}$ 用 $\delta_Q\dd\phi^I=\dd\zeta^I$ 求变分，$\dd\phi^{I_1}\wedge\cdots$ 中把每个 $\dd\phi^{I_i}$ 依次换成 $\dd\zeta^{I_i}$ 时 $\zeta$ 的反对易性自动实现反对称化，与 $f$ 全反对称（因 $\zeta^{I_1}\cdots\zeta^{I_k}$ 自动反对称化）一致，故左端的 $\delta_Q$ 恰等于右端的外微分。]

\section{带规范场的 4d A-模型}

正如章引言所指出的，我们希望通过对 $X$ 选取一个李群 $G$ 的作用、并把它与一个 4d 拓扑规范理论耦合，来对靶空间为 $X$ 的 4d A-模型做规范化。

一般而言，规范化过程并不唯一。我们对规范化的选择，是由希望与 Vafa-Witten 理论建立联系这一目标所引导的；也就是说，我们希望把 Vafa-Witten 理论描述为一类带规范场 4d A-模型的成员。为此，我们选择把 4d A-模型与一个规范群为 $G$ 的 Donaldson-Witten 拓扑规范理论耦合，并且在写出理论的作用量时，还纳入某些“非极小”的 BRST-正合项。

% [原文 p.25]
\subsection{Donaldson-Witten 拓扑规范理论}

我们回顾规范群为 $G$ 的 Donaldson-Witten 拓扑规范理论 [45] 的场含量、变分与作用量；它是 4d 中 $\cN=2$ SYM 的唯一扭转。

玻色场由规范场 $A$ 以及如下取值于 $\mathit{ad}\,P$ 的形式组成：
\begin{gather*}
  \sigma\in\Omega^0(M,\mathit{ad}\,P)\otimes\C,\\
  H\in\Omega^{2+}(M,\mathit{ad}\,P),
\end{gather*}
其中“$2+$”表示 $H$ 是关于 $M$ 上度量 $h$ 的自对偶 2-形式，而复 0-形式 $\sigma$ 由一对实 0-形式组成。费米子由
\begin{gather*}
  \eta\in\Omega^0(M,\mathit{ad}\,P),\\
  \psi\in\Omega^1(M,\mathit{ad}\,P),\\
  \chi\in\Omega^{2+}(M,\mathit{ad}\,P)
\end{gather*}
组成。场 $A$ 与 $H$ 的鬼数为 0，$\sigma$ 的鬼数为 2，$\psi$ 的鬼数为 1，而 $\eta$ 与 $\chi$ 的鬼数均为 $-1$。BRST 变分为
\begin{align}
  \delta_QA &= \psi, & \delta_Q\psi &= -\dd_A\sigma, \notag\\
  \delta_Q\sigma &= 0, \notag\\
  \delta_Q\bar\sigma &= \eta, & \delta_Q\eta &= [\sigma,\bar\sigma], \tag{3.2}\\
  \delta_Q\chi &= H, & \delta_QH &= [\sigma,\chi] \notag
\end{align}
其中 $\bar\sigma=-\sigma^\dagger$。与 4d A-模型的变分不同，这些变分不是幂零的；相反，$\delta_Q^2=\delta_\sigma$ 作用于所有场。\footnote{由于 $\sigma$ 是复的，$\delta_\sigma$ 实际上是无穷小 $G_\C$ 变换，其中 $G_\C$ 是复化李群。作用量将关于 $G_\C$ 变换不变。}

% [Tutor Note: p.25 (3.2) 中 $\delta_Q^2=\delta_\sigma$ 可直接逐场核对：$\delta_Q^2A=\delta_Q\psi=-\dd_A\sigma=\delta_\sigma A$（$\delta_\sigma A=-\dd_A\sigma$），$\delta_Q^2\sigma=0$（$\sigma$ 是 $G_\C$ 标量，$\delta_\sigma\sigma=0$）；关键是 $\delta_Q^2$ 不是零而是“模规范变换意义下为零”，这正是原文脚注强调 $G_\C$ 的原因。]

BRST-正合作用量为
\begin{equation}
  S=\frac{1}{e^2}\delta_Q\int_M\dd^4x\sqrt{h}\,\Tr\left\{\chi_{\mu\nu}\left(H^{\mu\nu}-2F^{+\mu\nu}\right)+2\bar\sigma D_\mu\psi^\mu\right\} \tag{3.3}
\end{equation}
其中 $F^+$ 是 $F$ 的自对偶部分。此外，我们在作用量中还纳入一项用以测量丛 $P$ 的拓扑：
\[
  -\frac{\ii\tau}{4\pi}\int_M\Tr F\wedge F.
\]

% [原文 p.26]
其中回顾耦合 $\tau$ 由角参数 $\theta$ 与耦合 $e$ 定义如下：
\[
  \tau=\frac{\theta}{2\pi}+\frac{4\pi\ii}{e^2}.
\]

\subsection{物质扇区}

选取李群 $G$ 在流形 $X$ 上的一个作用：
\[
  l:G\times X\to X.
\]
对固定的 $x\in X$，我们有映射 $l(-,x):G\to X$，它把 $G$ 的单位元送到 $x$。该映射在单位元处的微分是自同态 $\fg\to T_xX$；等价地，它定义了元素 $V\in T_xX\otimes\fg^*$。我们把这个元素写成
\[
  V=V^I_a\frac{\partial}{\partial\phi^I}\tilde{T}^a,
\]
其中 $\tilde{T}^a$ 是 $\fg^*$ 的基，与 $\fg$ 的基 $T_a$ 对偶。向量场 $V_a$ 的李括号为
\[
  [V_a,V_b]=-f^c{}_{ab}V_c.
\]

对固定的 $g\in G$，我们有映射 $l_g\equiv l(g,-):X\to X$，可以认为它通过复合作用在 $\sigma$-模型映射上（即 $\Phi\to\Phi'=l_g\circ\Phi$）。该映射在 $x$ 处的微分是自同态 $\dd l_g:T_xX\to T_{l_g(x)}X$，可以认为它作用在取值于 $\Phi^*TX$ 的场 $\zeta$ 上为
\[
  \zeta'^I=\frac{\partial\phi'^I}{\partial\phi^J}\zeta^J.
\]
这些变换在场上的无穷小形式为
\begin{equation}
  \delta_\epsilon\phi^I=\epsilon^aV^I_a,\qquad
  \delta_\epsilon\zeta^I=\epsilon^a\zeta^J\partial_JV^I_a, \tag{3.4}
\end{equation}
其中 $\epsilon\in\fg$ 是 $G$-作用的参数。其余 $\sigma$-模型场取值于 $\Phi^*TX$，因此服从与上面 $\zeta$ 相同的变换律。

我们取 $l$ 通过等距作用在 $X$ 上，于是对固定的参数 $\epsilon\in\fg$，上述变分 $\delta_\epsilon$ 是 4d A-模型作用量的整体对称性。为了把它提升为局域对称性（$\epsilon(x)$ 在 $M$ 上变化），我们如下使 4d A-模型作用量中出现的导数关于 (3.4) 规范协变。我们把场 $\Phi:M\to X$ 替换为场 $\Phi:M\to E$，它不再只是一个到 $M$ 的映射，而是相伴纤维丛 $P\times_GX$ 的一个截面（其典型纤维为通过作用 $l$ 与 $P$ 相伴的 $X$）。这里 $E$ 是全空间。相伴纤维丛的详细讨论见附录 C。

% [Tutor Note: p.26 “no longer merely a map into M”应为“map into X”。判断依据：$E\to M$ 是以 $X$ 为典型纤维的丛，$E$ 的截面是映射 $M\to E$；若写成“到 $M$ 的映射”则与“截面”矛盾（截面不可能落入底空间 $M$）。原文系笔误，译文照原文字面译出并在此指出。]

% [原文 p.27]
同样，我们把取值于拉回丛 $\Phi^*TX$ 的场 $\zeta$ 替换为取值于拉回丛 $\Phi^*VE$ 的场 $\zeta$，其中（如附录 C 所定义）垂直丛 $VE$ 由 $E$ 的切向量中沿纤维方向的那些组成。对其余 $\sigma$-模型场做类似的替换。

底主丛 $P$ 在邻域 $U^{(\alpha)}\subset M$ 上的平凡性诱导出映射 $x^{(\alpha)}:E^{(\alpha)}\to X$，其中 $E^{(\alpha)}\equiv\pi_E^{-1}(U^{(\alpha)})$。局部定义的场 $\Phi^{(\alpha)}\equiv x^{(\alpha)}\circ\Phi$ 是映射
\[
  \Phi^{(\alpha)}:U^{(\alpha)}\to X,
\]
而局部定义的场 $\zeta^{(\alpha)}\equiv\dd x^{(\alpha)}\circ\zeta$ 是截面
\[
  \zeta^{(\alpha)}\in\Omega^0\left(U^{(\alpha)},\Phi^{(\alpha)*}TX\right),
\]
其余场类似。利用靶空间上的局部坐标，我们可以进一步把场取值所在的丛平凡化：若 $V^{(\gamma)}\subset X$ 是带坐标 $\phi^{(\gamma)I}$ 的坐标卡，则在 $U^{(\alpha,\gamma)}\equiv U^{(\alpha)}\cap[\Phi^{(\alpha)}]^{-1}(V^{(\gamma)})$ 上，我们把场 $\zeta^{(\alpha)}$ 按分量写成
\[
  \zeta^{(\alpha)}=\zeta^{(\alpha,\gamma)I}\frac{\partial}{\partial\phi^{(\gamma)I}}.
\]

如附录 C 所述，联络 $A$ 给出垂直投影 $A_E^V:TE\to VE$ 的概念，进而允许我们定义作用在 $E$ 的截面上的协变导数如下：
\[
  \dd_A\Phi\equiv A_E^V\circ\dd\Phi\in\Omega^1\left(\Sigma,\Phi^*(VE)\right).
\]
局部地，协变导数由取值于切丛的 1-形式表示
\[
  \dd_A\Phi^{(\alpha)}\equiv\dd x^{(\alpha)}\circ\dd_A\Phi\in\Omega^1\left(U^{(\alpha)},\Phi^{(\alpha)*}TX\right),
\]
或更局部地，由普通 1-形式表示
\[
  \mathcal{D}\phi^I_{(\alpha,\gamma)}\equiv\left(\dd_A\Phi^{(\alpha)}\right)^*\dd\phi^{(\gamma)I}\in\Omega^1\left(U^{(\alpha,\gamma)}\right).
\]
$\Phi^*VE$ 的截面的协变导数类似地定义。局部 1-形式具体给出为

% [Tutor Note: p.27 这里 $\dd_A\Phi\in\Omega^1(\Sigma,\Phi^*(VE))$ 中的 $\Sigma$ 应为 $M$：本节基地流形始终记作 $M$（见 p.21“base manifold $M$”与 (3.5) 之后的“only defined locally on $M$”），$\Sigma$ 是附录 C 中一般流形的记号。判断依据：同一段下一行的局部形式明确写在 $U^{(\alpha)}\subset M$ 上。译文照原样保留 $\Sigma$ 并在此指出。]

% [Tutor Note: p.27 局部截断 $x^{(\alpha)}:E^{(\alpha)}\to X$ 的存在性依据是主丛的局部平凡化：$P|_{U^{(\alpha)}}\cong U^{(\alpha)}\times G$，于是 $E^{(\alpha)}=P\times_GX|_{U^{(\alpha)}}\cong U^{(\alpha)}\times X$，第二投影即 $x^{(\alpha)}$。这也解释了为什么下文的 $\Phi^{(\alpha)}$ 只是局部定义的。]
